\section{An Abstract Presentation of Mining}\label{sec:mining}

Different mining algorithms fit
inside the uniform setting of a competitive game
played, concurrently and independently, by the peers of a blockchain network, whose
goal is to decide which peer may add a \emph{valid} \textbf{next}
block to the stable chain. \emph{Valid} means that it respects
all consensus rules. For instance,
it points back to an existing block and
includes transactions with valid signatures
(the full, large list of consensus rules is outside the topic of this paper).
The game works in terms of:
%
\begin{itemize}
\item \emph{challenge}: from the current chain $\xi$,
  compute, deterministically, a challenge for \textbf{next};
\item \emph{answer}: create, if possible,
  a \textbf{next} that satisfies the challenge;
\item \emph{competition}: among more possible \textbf{next}s, prefer the \emph{best}
  one \wrt{} some notion of chain quality.
\end{itemize}
%
Each peer runs this mining algorithm,
whispers its resulting \textbf{next} to its
connected peers and receives the \textbf{next} of its connected peers. In case
of more alternative \textbf{next}s, the competition rule lets one choose the
winner. This gives rise to a \emph{consensus algorithm} that uses the mining
algorithm as a subroutine.

For proof of work, the mining game becomes:
%
\begin{itemize}
\item \emph{challenge}: compute an integer $\difficulty$ from $\xi$
  and find a valid \textbf{next} such that $\hash(\mathbf{next})\ge\difficulty$;
\item \emph{answer}: create a \textbf{next} by rotating the possible
  values for its nonce field, until $\hash(\mathbf{next})\ge\difficulty$;
\item \emph{competition}: among more \textbf{next}s, prefer that with the largest hash.
\end{itemize}
%
The computation of $\difficulty$ is not relevant here.
Moreover, Bitcoin's proof of work is usually given
in terms of a \emph{smallest} hash, but we require (equivalently) the
\emph{largest} hash to match the intuition that a higher $\difficulty$ makes mining harder.

For proof of stake, it is generally accepted that there is
\emph{no} mining~\cite{Kwon14}, but a mining game can be defined also in this case:
%
\begin{itemize}
\item \emph{challenge}: compute an integer $i$ from $\xi$ (the index of the chosen
  validator);
\item \emph{answer}: the $i$th validator
  creates a valid \textbf{next}; the others do not create anything and
  wait to receive \textbf{next};
\item \emph{competition}: among more possible \textbf{next}s, choose any of them, since
  this should not occur for this mining algorithm.
\end{itemize}
%
The algorithm for computing $i$ is irrelevat here: often, it is a simple
round-robin. Moreover, the consensus algorithm typically enriches the above
mining algorithm with a Byzantine protocol of pre-votes and pre-commits~\cite{Kwon14}.

Finally, proof of space, as presented in~\cite{Spoto25},
is based on the idea that \textbf{next} must contain
a data structure, called \emph{nonce}, that can be quickly recovered from a
set of precomputed nonces on disk (the \emph{plot file}). The mining game in this case is:
%
\begin{itemize}
\item \emph{challenge}: compute some data $\kappa$ from $\xi$;
\item \emph{answer}: create a valid \textbf{next} block, containing a nonce $\nu$
  derived from the plot file and from $\kappa$;
\item \emph{competition}: among more possible \textbf{next}s, choose that with the smallest
  $\delay(\nu,\kappa)$, where $\nu$ is the nonce in \textbf{next}.
\end{itemize}
%
The challenge $\kappa$ is used to select a nonce, from the plot file, that minimizes
$\delay$. The key point of this mining algorithm
is that the computation of the nonces (\ie{}, of the plot file) is so expensive
that it cannot be done on-the-fly, for each mined block, but rather
once and for all, before starting mining. The plot remains on disk to answer, quickly, mining
challenges. The consensus algorithm uses $\delay$ to wait until the block becomes valid.

This abstract presentation of mining clarifies
that mining is a subroutine whose input is the challenge derived from the current
chain $\xi$ and whose output is the answer \ie{}, the information
needed to create \textbf{next}:
a nonce for proof of work and proof of space (although a different notion of nonce is
used in either case), and a Boolean for proof of stake
(create the block or rather let another validator do it).
Fig.~\ref{fig:mining} shows this:
the miner that provides the answer to create \textbf{next}
and the peer that actually creates \textbf{next}
are separate software parts of the same tool. Traditionally,
such software parts work in the same machine.
But they have been already moved to different machine, for instance in pool
mining architectures, where a single peer leverages the power of many, remote
miners, working in different machines and connected through a network
(Fig.~\ref{fig:mining}).
Namely, many users want to run a miner and earn cryptos,
but only a few are interested in running a full blockchain peer on a dedicated machine,
which requires technical skills, database space and a very good internet connection.

\begin{figure}[t]
  \begin{center}
    \begin{tikzpicture}
      \draw [fill=verylightviolet,thick] (-1,-0.3) rectangle +(4,4.7);
      \draw (1,2.5) node {blockchain peer};
      \draw (1,1.9) node {{\small (networking, consensus}};
      \draw (1,1.5) node {{\small and application)}};

      \draw [fill=verylightred,thick] (7,3) rectangle +(2,0.8);
      \draw (8.0,3.4) node {$\text{miner}_1$};
      \draw[->, thick] (3,3.65) -- (7,3.65);
      \draw (5.25,3.85) node {{\small challenges}};
      \draw[<-, thick] (3,3.15) -- (7,3.15);
      \draw (5.25,2.9) node {{\small answers}};
      \draw (5.25,2.65) node {{\scriptsize (in this paper, nonces or deadlines)}};

      \draw (8.0,2.075) node {$\vdots$};

      \draw [fill=verylightred,thick] (7,0.35) rectangle +(2,0.8);
      \draw (8.0,0.75) node {$\text{miner}_n$};
      \draw[->, thick] (3,1.0) -- (7,1.0);
      \draw (5.25,1.2) node {{\small challenges}};
      \draw[<-, thick] (3,0.5) -- (7,0.5);
      \draw (5.25,0.25) node {{\small answers}};
      \draw (5.25,0) node {{\scriptsize (in this paper, nonces or deadlines)}};
    \end{tikzpicture}
  \end{center}
  \caption{Blockchain peer and miner(s) as separate software tools, possibly
    running in different machines connected through a network.}\label{fig:mining}
\end{figure}

Up to now, the machines where the miners are running
are computers or specialized, very expensive hardware. The goal of this article
is to allow such miners to run inside a mobile device as well.
But this might result impossible or impractical,
depending on the mining algorithm and on the machine
where the miner will run. Let us consider, in particular, the case when
the mining algorithm is one of those discussed
in Sec.~\ref{sec:introduction}.

\subsubsection*{Proof of work}
%
In this case, the challenge includes the full block being
created, since its nonce must maximize its hash.
For Bitcoin, a block's size is up to 4~MB after the Segregated Witness upgrade.
This data must be exchanged at \emph{each} block creation, impossible in many
countries where mobile internet is slow; moreover, most mobile operators limit
the exchanged GB per month; finally, the computational cost
of proof of work would drain the battery of the mobile device. Consequently,
at the moment, mining with proof of work is irrealistic on a mobile device,
for a global blockchain such as Bitcoin.

\subsubsection*{Proof of stake}
%
In this case, mining is very cheap and limited to the validators. The exchanged data is minimal.
But validators have no reason to offload their very simple mining activity.
Therefore, mining on a mobile device is possible for proof of stake, but it does
not seem of any practical interest.

\subsubsection*{Proof of space}
%
In this last case, mining requires the use of disk space, but no significant
CPU power. Currently, mobile devices have an average disk space
of 256~GB, smaller than the average 1000~GB of desktop computers but still significant.
Moreover, the proof of space of~\cite{Spoto25} uses very small challenges (around 64~bytes)
and answers (nonces, compacted into \emph{deadlines} of around 200~bytes).
The block being created is not needed for mining. This makes that
algorithm ideal for a remote miner in a mobile device:
no need of significant CPU power, enough disk space in the device and little data exchange.
This is why it has been the choice for the implementation of our Mokaminter app.
